\section{Formal Reconstruction and Acceptance Method}
\label{sec:formal-method}

\subsection{Typed partial reconstruction}

All byte objects in the model are finite octet strings in
$\mathbb{B}^{*}$, where $\mathbb{B}=\{0,\ldots,255\}$. Length $|x|$ is measured
in octets. For each admitted profile $k\in\mathcal{K}_p$, reconstruction is a
partial, resource-bounded computation
\begin{equation}
  R_{p,k}(x)\Downarrow c
  \quad\text{or}\quad
  R_{p,k}(x)\Uparrow e,
  \label{eq:reconstruction-relation}
\end{equation}
where $c$ is the typed candidate from \cref{eq:candidate-artifact} and $e$ is a
typed terminal error. The computation can reject, time out, be cancelled, or
fail during bounded external processing. It is not a total function from
arbitrary bytes to valid media.

The performed reconstruction level is ordered by
\begin{equation}
  \Unprocessed \prec \Structural \prec \Semantic.
  \label{eq:reconstruction-order}
\end{equation}
At the structural level, a handler parses the source container, accounts for
the ranges relevant to its transformation, emits only admitted components,
rewrites dependent lengths or offsets, and excludes policy-forbidden metadata.
Encoded media payload can remain byte-identical. At the semantic level, a
bounded decoder produces an intermediate media state $s$, and a pinned encoder
constructs a new artifact from $s$ and policy constants. The model makes no
claim that either level establishes benign media meaning.

The policy requirement is
$q_p(k)\in\{\Structural,\Semantic\}$. The second value represents
semantic-or-block behavior. For the level $c.\ell$ reported by a candidate,
\begin{equation}
  \mathsf{LevelOK}_p(k,c.\ell)
  \iff q_p(k)\preceq c.\ell.
  \label{eq:level-satisfaction}
\end{equation}
Thus semantic work satisfies a structural minimum, but structural work never
satisfies a semantic requirement. Failure or absence of the semantic path
cannot silently select a structural fallback.

\subsection{Policy-indexed canonical output language}

The complete parser $\mathsf{Parse}_{v,k}$ belongs to the fixed registry version
$v$ and yields the report in \cref{eq:parse-report}. Policy $p$ is then applied
to that stable report. The canonical output language is
\begin{align}
  L_{p,k}=\{y\in\mathbb{B}^{*}:{}&
  \mathsf{Parse}_{v,k}(y)\Downarrow\rho_k(y),\notag\\
  &\rho_k.o=0,\ \rho_k.t=|y|,\notag\\
  &\mathsf{Allowed}_{p,k}(\rho_k.C),\notag\\
  &\mathsf{MetadataOK}_{p,k}(\rho_k.C),\notag\\
  &\mathsf{ParseBounds}_{p,k}(\rho_k.q)\}.
  \label{eq:canonical-language}
\end{align}
The language is defined by admitted structure, not by subtracting a short list
of known malicious features. Unknown components, invalid locations, duplicate
mandatory elements, stale references, arithmetic overflow, unconsumed ranges,
and exhausted parser budgets are excluded unless reconstruction removes them
before the final parse.

Metadata exclusion is structural. For each profile and component location,
policy defines $\mathcal{M}^{\mathrm{forbid}}_{p,k}$.
$\mathsf{MetadataOK}_{p,k}$ requires the parsed component tree to contain none
of these members. It does not assert that the byte spelling of a tag cannot
appear accidentally inside compressed media. Signature rules likewise have an
explicit scope such as start of file, end of file, parsed component location,
or an exact administrative indicator.

\begin{figure}[t]
\centering
\resizebox{0.96\textwidth}{!}{%
\begin{tikzpicture}[
  node distance=6mm and 7mm,
  box/.style={draw,rounded corners=2pt,align=center,minimum height=9mm,minimum width=25mm,font=\small},
  guard/.style={box,fill=blue!7},
  transform/.style={box,fill=green!8},
  terminal/.style={box,fill=gray!12},
  arrow/.style={-{Latex[length=2mm]},thick}
]
\node[guard] (route) {three-source\\route};
\node[transform,right=of route] (rebuild) {required-level\\reconstruction};
\node[guard,right=of rebuild] (parse) {independent\\complete parse};
\node[guard,right=of parse] (gates) {type, signature,\\and resource gates};
\node[terminal,right=of gates] (publish) {clean-only\\result and publication};
\draw[arrow] (route) -- (rebuild);
\draw[arrow] (rebuild) -- (parse);
\draw[arrow] (parse) -- (gates);
\draw[arrow] (gates) -- (publish);
\node[terminal,below=8mm of parse] (reject) {non-clean result or typed error\\no deliverable bytes};
\draw[arrow] (route) -- (reject);
\draw[arrow] (rebuild) -- (reject);
\draw[arrow] (parse) -- (reject);
\draw[arrow] (gates) -- (reject);
\end{tikzpicture}%
}
\caption{Conjunctive acceptance method. Each downward edge represents failure
or indeterminacy. No stage can authorize publication independently.}
\label{fig:formal-acceptance-path}
\end{figure}

\subsection{Acceptance predicate}

Let $c=(y,n_{\mathrm{out}},m_{\mathrm{out}},k_h,\ell,\eta)$. For the unique
route $k$, candidate acceptance is
\begin{align}
\mathsf{Acc}_{p}(a,c,k)\iff{}&
 \mathsf{Route}_{p}(a)=k \notag\\
&\land \mathsf{SourceResolved}_{p}(a,k,\eta) \notag\\
&\land \mathsf{LevelOK}_{p}(k,\ell) \notag\\
&\land \mathsf{Full}_{v}(y)=\{k\} \notag\\
&\land y\in L_{p,k} \notag\\
&\land \mathsf{TypeOK}_{p}(n_{\mathrm{out}},m_{\mathrm{out}},k_h,k,y) \notag\\
&\land \mathsf{SignatureOK}_{p}(y) \notag\\
&\land \mathsf{BoundsOK}_{p}(x,y,\eta) \notag\\
&\land \mathsf{TerminalOK}_{p}(\eta).
\label{eq:extended-acceptance}
\end{align}
$\mathsf{TypeOK}_p$ requires the generated suffix, declared output type,
handler profile, byte recognition, and complete parser profile to converge.
$\mathsf{BoundsOK}_p$ contains both pre-reconstruction and output limits.
$\mathsf{TerminalOK}_p$ excludes timeout, cancellation, panic, parser
exhaustion, child-process failure, warning-only behavior, and dry-run behavior.

The scientific accepted set is
\begin{equation}
  \mathcal{A}_{p}=
  \{(a,c):\exists k\ .\
  \mathsf{Route}_{p}(a)=k\land
  R_{p,k}(x)\Downarrow c\land
  \mathsf{Acc}_{p}(a,c,k)\}.
  \label{eq:accepted-set}
\end{equation}
The implementation must not try handlers until one happens to accept. A route
exists only for a singleton eligible set, and output acceptance is evaluated
again on the candidate produced by that route.

\subsection{Formal invariants F1--F8}

The acceptance relation is organized around eight design invariants. These
invariants state the obligations of a clean result. They do not claim that the
deferred empirical campaigns have passed.

\begin{table}[t]
\centering
\caption{Formal invariant register and corresponding evidence obligation.}
\label{tab:properties-p1-p8}
\small
\begin{tabularx}{\textwidth}{L{0.07\textwidth}L{0.25\textwidth}YL{0.22\textwidth}}
\toprule
ID & Invariant & Formal obligation & Principal evidence \\
\midrule
F1 & Complete consumption & $\rho_k.o=0\land\rho_k.t=|y|$ & Parser report and suffix mutation \\
F2 & Canonical component closure & $\mathsf{Allowed}_{p,k}(\rho_k.C)$ & Location-aware parser and component mutation \\
F3 & Type convergence & $\mathsf{Full}_{v}(y)=\{k\}$ and $\mathsf{TypeOK}_{p}$ & Alias parity and complete output parsing \\
F4 & Metadata exclusion & $\rho_k.C\cap\mathcal{M}^{\mathrm{forbid}}_{p,k}=\varnothing$ & Parsed tree and field mutation \\
F5 & Source resolution and provenance & $\mathsf{SourceResolved}_{p}(a,k,\eta)$ and $R_{p,k}(x)\Downarrow c$ & Source accounting and reconstruction trace \\
F6 & Resource boundedness & Every applicable counter satisfies its pre-operation bound & Static review, fault injection, and measurement \\
F7 & Independent candidate validation & Complete parser, type, and signature gates evaluate the produced $y$ & Direct validation, fuzzing, and differential tools \\
F8 & Fail-closed result construction & Acceptance is the only route to a clean result with deliverable bytes & Interface invariants and publication fault campaign \\
\bottomrule
\end{tabularx}
\end{table}

F1 rules out accepting a valid prefix while leaving an unclassified suffix. F2
requires each emitted component to be admitted at its parsed location and each
cross-reference to resolve within an allowed extent. F3 prevents one recognizer
from overruling incompatible output evidence. F4 applies to parsed metadata
classes rather than incidental substrings in compressed payload.

F5 makes the source-side obligation explicit. Structural copying requires a
complete account of copied and discarded ranges. Semantic reconstruction
requires a bounded decode of the selected interpretation and records that no
source container range is intentionally copied. F6 covers input and output
size, dimensions, aggregate pixels, component and frame counts, nesting,
duration, channels, sample tables, expansion, cancellation, and deadline where
applicable. Exhaustion is terminal. It cannot be reinterpreted as end of input.

F7 requires a post-reconstruction path that does not trust handler success or a
handler-provided media type. Independence in the formal model means a separate
validator over the produced candidate. Empirical corroboration uses separately
implemented parser or decoder families. F8 connects acceptance to the result
and publication interfaces defined in \cref{sec:publication-semantics}.

Second-pass reconstruction stability is deliberately not an acceptance
conjunct. For deterministic structural and lossless profiles, byte equality
under a second pass is an empirical obligation. For lossy semantic profiles,
the separately frozen comparator evaluates canonical closure and bounded
quality drift. These outcomes have different meanings and are never pooled as
``idempotence.''

\subsection{Conditional analytical consequences}

\begin{property}[Canonical membership of clean candidates]
\label{prop:accepted-canonical}
If $(a,c)\in\mathcal{A}_{p}$, then $c.y\in L_{p,k}$ for one observed profile
$k\in\mathcal{K}_p$. This follows directly from the acceptance conjunction and
is therefore a design invariant, not an independently proved security theorem.
Its force depends on parser soundness and on the completeness of the modeled
observation registry.
\end{property}

\begin{proposition}[Syntactic-source independence under deterministic semantic reconstruction]
\label{prop:semantic-source-independence}
Assume a deterministic bounded decoder $\delta_{p,k}$ and encoder
$\gamma_{p,k}$ with
$R_{p,k}(x)=\gamma_{p,k}(\delta_{p,k}(x),p)$. If
$\delta_{p,k}(x_1)=\delta_{p,k}(x_2)=s$, then the produced byte component of the
two candidates is equal.
\end{proposition}

\begin{proof}[Analytical argument]
Substitution gives
$\gamma_{p,k}(\delta_{p,k}(x_1),p)=\gamma_{p,k}(s,p)=
\gamma_{p,k}(\delta_{p,k}(x_2),p)$. Source-container differences that do not
alter decoded state therefore cannot influence a deterministic output encoder.
The argument does not establish harmlessness of $s$, decoder correctness, or
absence of coincidental byte patterns.
\end{proof}

\begin{proposition}[Bounded rejection under covered counters]
\label{prop:bounded-rejection}
Suppose every loop and allocation reachable before acceptance is governed by a
finite policy counter or enforced deadline, checked before the bounded
operation, and exhaustion is terminal. An input cannot be accepted by causing a
covered counter to exceed its limit.
\end{proposition}

\begin{proof}[Analytical argument]
The first operation that would exceed a covered counter is not executed.
Deadline or cancellation expiration likewise produces a terminal error at the
next enforced boundary. At least one of $\mathsf{BoundsOK}_p$ and
$\mathsf{TerminalOK}_p$ is false, so $\mathsf{Acc}_p$ is false. The proposition
does not establish coverage of every implementation operation. Coverage remains
a static, fuzzing, fault-injection, and measurement obligation.
\end{proof}

The formal consequences delimit the contribution. They do not estimate false
blocking, carrier-property retention, parser defect probability, latency,
memory demand, or implementation exploitability. Those questions belong to the
pre-specified empirical protocol.
